\section{第三范式 Online Learning：为什么它可能比机器人、世界模型、多模态更根本}

前面讨论的是现有竞争格局，本章进入技术范式。广密认为 Pre-training scaling 快结束了，RL/Post-training 仍有空间，但下一阶段真正可能改变格局的是 Online Learning 或 Continuous Learning。这里的核心不是“模型有更长上下文”这么简单，而是模型能否从交互、反馈和真实任务中持续学习。

\begin{figure}[H]
\centering
\includegraphics[width=0.96\textwidth]{scaling-paradigm-shift.png}
\caption{三种学习范式：Pre-training、RL/Post-training、Memory 与 Online Learning。自制概念图，依据 00:36:01--00:41:01 对谈内容整理。}
\end{figure}

\begin{knowledgebox}{读图：从静态压缩走向交互学习}
Pre-training 像一次性压缩大规模互联网和专业数据；RL/Post-training 用人类反馈、专家过程和任务结果塑形；Memory 和长上下文让模型记住更多用户与任务历史；Online Learning 进一步要求模型在推理和交互过程中真正更新自己，越用越懂用户、越做越会做。
\end{knowledgebox}

\subsection{Pre-training：巨大的压缩器}

Pre-training，即预训练，是用海量文本、代码、多模态数据训练模型，让它学会预测下一个 token 或下一段结构。广密强调，今天的模型仍然像巨大的压缩器：数据分布里没有的任务往往不 work，只有压缩过类似数据，模型才更可靠。

这句话很关键。它解释了为什么模型在知识层面已经超过多数人，却在真实工作流里仍会失败。因为真实工作不是只回答问题，而是要进入 CRM、财务系统、设计软件、医院流程、销售现场、企业权限和具体工具链。模型如果没有见过足够多真实过程，就只能靠泛化猜测。

\begin{importantbox}{压缩器视角}
模型不是凭空“懂”任务，而是把训练数据中的模式压缩成参数和表示。一个任务如果长期没有进入数据分布，模型就很难稳定完成；要让 Agent 可靠，必须把真实工作过程变成可学习的数据。
\end{importantbox}

\subsection{RL/Post-training：专家数据与任务反馈}

前面把预训练解释为大规模压缩，本节进一步看“压缩之后怎样塑形”。如果说 Pre-training 让模型学到世界上已经存在的大量模式，那么 RL/Post-training 更关心模型怎样在具体任务、用户偏好和专家过程里变得可用。

RL 是 Reinforcement Learning，强化学习。它不只是“奖励模型”四个字，而是一套让系统根据结果反馈调整行为的训练方式。在大模型语境下，Post-training 通常包括指令微调、人类偏好反馈、专家轨迹、自动评测、工具使用反馈和任务结果评分。节目中把这类数据比作“新能源”：有用，但总量较少，采集成本高。

如果 Coding Agent 要变强，需要高质量代码、真实 issue、测试、PR 反馈和长期软件工程过程；如果医疗 Agent 要变强，需要医生诊断过程、病例、检查、用药、随访和责任边界；如果投研 Agent 要变强，需要分析师工作流、数据来源、模型假设和错误复盘。每一类都不是普通网页文本，而是专家过程数据。

\begin{knowledgebox}{术语消化：RL/Post-training 相关词}
\begin{tabularx}{\textwidth}{p{0.22\textwidth}p{0.30\textwidth}X}
\toprule
术语 & 核心机制 & 本期作用 \\
\midrule
RL & 用奖励或结果反馈优化策略 & 让模型不只模仿文本，还能朝任务成功优化。 \\
Post-training & 预训练后的对齐、指令、偏好和任务训练 & 决定模型是否符合产品目标。 \\
专家轨迹 & 金牌专家完成任务的过程记录 & 让模型学习真实工作流，而不只是答案。 \\
Credit assignment & 把成功或失败归因到中间步骤 & Robotics 章节中用来解释每一步是否有利于任务成功。 \\
Value function & 估计当前状态或动作未来价值的函数 & 帮助系统知道哪些步骤更接近成功。 \\
\bottomrule
\end{tabularx}
\end{knowledgebox}

\subsection{Online Learning：边交互边学习}

Online Learning 在本期中指模型自主学习、持续学习、从交互中吸收用户和任务经验。与静态预训练不同，它要求模型在使用过程中变得更懂用户、更懂任务、更主动。广密举的产品图景是：今天 ChatGPT 必须等用户 prompt；未来 AI 可能在后台反思、准备、联想，下一次用户提问时已经有了更深的理解。

这个能力一旦成立，会改变产品形态。AI 不再是每次从零推理的工具，而是和用户共同成长的长期伙伴；不再只是回答当前问题，而是记住用户目标、偏好、项目、上下文和历史错误；不再只是被动等命令，而是主动准备和提醒。

\begin{importantbox}{Online Learning 的定义}
在本讲义中，Online Learning 指模型或 AI 系统能从持续交互、任务结果和用户反馈中学习，并把学习结果稳定用于未来行为。它不是简单保存聊天记录，而是把经验转化为能力、偏好或策略更新。
\end{importantbox}

\begin{warningbox}{Memory 不等于 Online Learning}
Memory 可以只是存储用户信息；长上下文可以只是把更多历史塞进 prompt。Online Learning 要求系统从这些历史里学到可泛化的行为改进。只会“记住事实”不等于会“越用越聪明”。
\end{warningbox}

\subsection{能源隐喻：石油、新能源与核聚变}

本节用能源隐喻把前三个学习阶段压缩成一个容易记住的比较框架。这个隐喻的价值不在于精确预测技术路线，而在于帮助读者同时看到存量数据、专家数据和持续学习三类资源的约束差异。

广密用三个能源隐喻解释范式差异：Pre-training 数据像石油，量大但有限；RL 专家数据像新能源，有用但总量少；Online Learning 像核聚变，还没有真正突破，但如果突破，能量密度和上限会完全不同。

\begin{figure}[H]
\centering
\includegraphics[width=0.96\textwidth]{energy-metaphor.png}
\caption{能源隐喻：石油、新能源和核聚变对应三类数据与学习方式。自制概念图，依据 00:41:01--00:43:05 对谈内容整理。}
\end{figure}

\begin{knowledgebox}{读图：隐喻的边界}
石油强调存量数据有限，新能源强调专家数据有用但采集困难，核聚变强调 Online Learning 的潜在上限。隐喻不说明技术一定会按能源史重演；它只是帮助我们区分三类数据和学习机制的约束。
\end{knowledgebox}

\subsection{为什么说机器人、世界模型和多模态可能是“假问题”}

节目中最容易引发争议的一句话是：机器人、世界模型、多模态很多可能是假问题，Online Learning 才是真问题。这里的“假问题”不是说这些方向不重要，而是说它们可能依赖更底层的泛化和持续学习。如果没有泛化，机器人就会像自动驾驶一样：每个长尾场景都要采数据、标数据、调规则，长期陷入“有多少人工就有多少智能”。

多模态和世界模型也类似。模型能看视频、生成图像、预测物理过程，不等于它能从真实环境中持续学习、迁移经验、处理新任务。真正难的是把经验变成可复用能力。Online Learning 如果突破，机器人、世界模型和多模态都可能获得更强泛化；如果不突破，这些方向就会继续被数据量、标注成本和真实环境复杂度拖住。

\subsection{本章小结}

第三范式的核心不是一个新名词，而是从静态压缩走向持续交互学习。Pre-training 解决大规模知识压缩，RL/Post-training 解决任务反馈和专家过程，Online Learning 则试图解决“越用越聪明”和“从真实环境中泛化”的问题。它之所以重要，是因为它可能改变现有三强交替领先的平衡。

